% Standalone one-page cover note: points readers of Part I to Part II.
% Does NOT modify the published Part I PDF; meant to be prepended / attached / used as a note.
\documentclass[11pt]{article}
\usepackage[T1]{fontenc}
\usepackage{lmodern}
\usepackage{amsmath,amssymb,graphicx,hyperref}
\usepackage{xurl}
\usepackage[table]{xcolor}
\definecolor{ggteal}{HTML}{1B6F8C}
\definecolor{ggband}{HTML}{DCE7EB}
\usepackage[margin=1.1in]{geometry}
\usepackage{microtype}
\usepackage[most]{tcolorbox}
\usepackage{parskip}
\pagestyle{empty}
\hypersetup{colorlinks=true,linkcolor=ggteal,urlcolor=ggteal}
\newcommand{\lean}[1]{\texttt{#1}}

\begin{document}

{\centering
{\large\bfseries A note to the reader of Part~I}\\[2pt]
{\color{ggteal}\rule{0.5\textwidth}{0.8pt}}\\[6pt]
{\itshape The sequel proposed at the end of this paper is now written.}\\[14pt]
}

Part~I (\emph{Random Signs into Matchings}) closed with a proposal and a promise:
to prove the \textbf{Heilmann--Lieb theorem} not through interlacing machinery but
through \textbf{Godsil's path tree}, and---\emph{``if this paper has a sequel, this
is its first sentence.''} That sequel now exists.

\textbf{Part~II --- \emph{Unfolding a Graph into a Tree: A Machine-Checked Proof of
the Heilmann--Lieb Theorem in Lean~4}} carries out the proposal in full. It builds
the path tree, proves the divisibility $\mu_G \mid \mu_{T(G,u)}$ (crossing the one
elaboration wall the path-tree route had stalled at), establishes the forest
identity, and proves the Ramanujan bound from a weighted Gershgorin estimate. Both
halves are \lean{sorry}-free.

\begin{tcolorbox}[colframe=ggteal,colback=ggband!40,arc=3pt,boxrule=0.9pt]
\centering
For every finite graph $G$, the matching polynomial $\mu_G$ is real-rooted; and if
$2\le\Delta$ and every vertex has degree $\le\Delta$, every root of $\mu_G$ lies in
\[
  \bigl[-2\sqrt{\Delta-1},\;2\sqrt{\Delta-1}\bigr].
\]
\end{tcolorbox}

This is exactly the endpoint that Part~I mapped as future work
(``The next stone''): the first machine-checked Heilmann--Lieb, completing the
path-tree program. As in Part~I, the mathematics is classical and no new theorem is
claimed; the contribution is the formalization.

\bigskip
{\color{ggteal}\rule{\textwidth}{0.4pt}}\\[4pt]
\textbf{Where to find Part~II.}
\begin{itemize}\setlength{\itemsep}{1pt}
\item Sources and both PDFs (English + Spanish):
  \url{https://github.com/karlesmarin/godsil-gutman-lean}
\item Archived on Zenodo: DOI \texttt{10.5281/zenodo.20561832}
  (\url{https://doi.org/10.5281/zenodo.20561832}).
\end{itemize}

\vfill
{\footnotesize\color{gray}
Carles Mar\'in, \texttt{karlesmarin@gmail.com}. This note accompanies, and does not
alter, the archived Part~I (Zenodo DOI \texttt{10.5281/zenodo.20517350}).}

\end{document}
